% Superposición de Kabsch en 2D con puntos de generar.py
\begin{tikzpicture}[font=\sffamily\small,
  pq/.style={circle,fill=azul,draw=white,minimum size=2.4mm,inner sep=0pt},
  pp/.style={circle,fill=naranja,draw=white,minimum size=2.4mm,inner sep=0pt},
  corr/.style={gris,densely dotted,thin}]
  \begin{scope}[shift={(0,0)},scale=0.62]
    \coordinate (q0) at (0.000,0.000);
\coordinate (q1) at (1.500,0.300);
\coordinate (q2) at (2.400,1.600);
\coordinate (q3) at (1.200,2.500);
\coordinate (q4) at (-0.300,1.400);
\coordinate (p0) at (5.603,-1.121);
\coordinate (p1) at (6.150,0.366);
\coordinate (p2) at (5.790,1.846);
\coordinate (p3) at (4.222,1.354);
\coordinate (p4) at (4.275,-0.385);
\coordinate (c0) at (1.355,-0.373);
\coordinate (c1) at (1.902,1.114);
\coordinate (c2) at (1.542,2.594);
\coordinate (c3) at (-0.026,2.102);
\coordinate (c4) at (0.027,0.363);
\coordinate (a0) at (-0.034,-0.071);
\coordinate (a1) at (1.485,0.376);
\coordinate (a2) at (2.458,1.548);
\coordinate (a3) at (1.128,2.513);
\coordinate (a4) at (-0.236,1.433);
\def\capquinceRantes{4.48}\def\capquinceRcentro{1.21}\def\capquinceRfinal{0.08}\def\capquinceAng{53}

    \draw[azul!60] (q0)--(q1)--(q2)--(q3)--(q4)--cycle;
    \draw[naranja!70] (p0)--(p1)--(p2)--(p3)--(p4)--cycle;
    \foreach \i in {0,...,4}{\draw[corr] (q\i) -- (p\i);}
    \foreach \i in {0,...,4}{\node[pq] at (q\i) {}; \node[pp] at (p\i) {};}
  \end{scope}
  \begin{scope}[shift={(6.0,0)},scale=0.85]
    \coordinate (q0) at (0.000,0.000);
\coordinate (q1) at (1.500,0.300);
\coordinate (q2) at (2.400,1.600);
\coordinate (q3) at (1.200,2.500);
\coordinate (q4) at (-0.300,1.400);
\coordinate (p0) at (5.603,-1.121);
\coordinate (p1) at (6.150,0.366);
\coordinate (p2) at (5.790,1.846);
\coordinate (p3) at (4.222,1.354);
\coordinate (p4) at (4.275,-0.385);
\coordinate (c0) at (1.355,-0.373);
\coordinate (c1) at (1.902,1.114);
\coordinate (c2) at (1.542,2.594);
\coordinate (c3) at (-0.026,2.102);
\coordinate (c4) at (0.027,0.363);
\coordinate (a0) at (-0.034,-0.071);
\coordinate (a1) at (1.485,0.376);
\coordinate (a2) at (2.458,1.548);
\coordinate (a3) at (1.128,2.513);
\coordinate (a4) at (-0.236,1.433);
\def\capquinceRantes{4.48}\def\capquinceRcentro{1.21}\def\capquinceRfinal{0.08}\def\capquinceAng{53}

    \draw[azul!60] (q0)--(q1)--(q2)--(q3)--(q4)--cycle;
    \draw[naranja!70] (c0)--(c1)--(c2)--(c3)--(c4)--cycle;
    \foreach \i in {0,...,4}{\draw[corr] (q\i) -- (c\i);}
    \foreach \i in {0,...,4}{\node[pq] at (q\i) {}; \node[pp] at (c\i) {};}
  \end{scope}
  \begin{scope}[shift={(10.4,0)},scale=0.85]
    \coordinate (q0) at (0.000,0.000);
\coordinate (q1) at (1.500,0.300);
\coordinate (q2) at (2.400,1.600);
\coordinate (q3) at (1.200,2.500);
\coordinate (q4) at (-0.300,1.400);
\coordinate (p0) at (5.603,-1.121);
\coordinate (p1) at (6.150,0.366);
\coordinate (p2) at (5.790,1.846);
\coordinate (p3) at (4.222,1.354);
\coordinate (p4) at (4.275,-0.385);
\coordinate (c0) at (1.355,-0.373);
\coordinate (c1) at (1.902,1.114);
\coordinate (c2) at (1.542,2.594);
\coordinate (c3) at (-0.026,2.102);
\coordinate (c4) at (0.027,0.363);
\coordinate (a0) at (-0.034,-0.071);
\coordinate (a1) at (1.485,0.376);
\coordinate (a2) at (2.458,1.548);
\coordinate (a3) at (1.128,2.513);
\coordinate (a4) at (-0.236,1.433);
\def\capquinceRantes{4.48}\def\capquinceRcentro{1.21}\def\capquinceRfinal{0.08}\def\capquinceAng{53}

    \draw[azul!60] (q0)--(q1)--(q2)--(q3)--(q4)--cycle;
    \draw[naranja!70] (a0)--(a1)--(a2)--(a3)--(a4)--cycle;
    \foreach \i in {0,...,4}{\draw[corr] (q\i) -- (a\i);}
    \foreach \i in {0,...,4}{\node[pq] at (q\i) {}; \node[pp] at (a\i) {};}
  \end{scope}
  \node[font=\sffamily\bfseries\small,text=tinta] at (2.1,3.3) {(a) posición inicial};
  \node[font=\sffamily\bfseries\small,text=tinta] at (6.9,3.3) {(b) centroides unidos};
  \node[font=\sffamily\bfseries\small,text=tinta] at (11.6,3.3) {(c) rotación óptima};
  \node[etiqueta] at (2.1,-1.25) {RMSD $=\SI{4,48}{\angstrom}$};
  \node[etiqueta] at (6.9,-1.25) {RMSD $=\SI{1,21}{\angstrom}$};
  \node[etiqueta] at (11.4,-1.25) {RMSD $=\SI{0,08}{\angstrom}$ (giro de \ang{-53,4})};
  \draw[flecha=tinta2] (4.35,1.0) -- node[above,etiqueta]{$-\bar{\mathbf p}+\bar{\mathbf q}$} (5.3,1.0);
  \draw[flecha=tinta2] (8.6,1.0) -- node[above,etiqueta]{$\mathbf R=\mathbf V\mathbf D\mathbf U^{\top}$} (9.6,1.0);
  \node[pq,label={[etiqueta]right:referencia $\mathbf q_k$}] at (3.2,-1.9) {};
  \node[pp,label={[etiqueta]right:estructura móvil $\mathbf p_k$}] at (6.2,-1.9) {};
\end{tikzpicture}
